\honest{Bayesians vs. Barbarians}{Eliezer Yudkowsky}{2009}

I open with a passage from my earlier post on armies, which ends: ``an uncoordinated mob gets slaughtered.''

Suppose Evil Barbarians attack a country of rationalists. One view says the rationalists will lose: the Barbarians expect a reward for dying bravely, while each rationalist would rather someone else fought, and would keep ``that sausage'' for themselves, since no one sausage decides the war. I cite no one who holds this view. \nb{Its core is the free-rider problem. Mancur Olson wrote in 1965 that without ``coercion or some other special device,'' rational, self-interested people ``will not act to achieve their common or group interests.''}

War is not fun, and losing is worse. If rationalists always lost wars, I could not advocate rationality for everyone. Wars have been won without torture, but real wars ``cannot be won by refined politeness.''

In principle: I one-box on Newcomb's problem. I do not think the voters in an election won 100,000 to 99,998 were all irrational to vote. Two AIs using ``my kind of decision theory'' that can read each other's code will cooperate in the true Prisoner's Dilemma, ``(Or even just common knowledge of each other's rationality in the appropriate sense.)'' \nb{For programs with identical code, mutual cooperation was already known to be an equilibrium. The claim about rationality alone is disputed: the \textsc{Stanford Encyclopedia} calls it ``controversial,'' and the other side holds that one player's choice cannot cause the other's.} Agents like that ``will end up at Pareto optima instead of Nash equilibria.'' \nb{The theory was not published until 2010; the reader cannot check this.} So a community of AIs could pick soldiers by lottery, having changed their code beforehand so that the chosen will fight.

For humans the pure dilemma is rare, because of reputation, caring and incentives. If children are not taught that rationalists defect, I think people would volunteer. If too few do, ``we can perhaps all agree,'' before the lottery, to give the selectees courage drugs and ``shoot them if they run away.'' \nb{That is Olson's ``coercion.'' Real armies have shot soldiers who ran, the British Army in the First World War among them.}

I do not endorse the draft as it is practiced: ``Drafts are a tool of kings playing games in need of toy soldiers.'' \nb{The Vietnam draft was itself a lottery from December 1969. What the post's scheme adds is the prior agreement of those drafted, and the post gives no evidence that present-day drafts lack it.} Soldiers must also obey orders, and if they will not, we may agree in advance to shoot those who disobey. ``This is not to say that those who fragged their own officers in Vietnam were in the wrong; for they could have consistently held that they preferred no one to participate in the draft lottery.''

I write all this for the sake of self-image: a society of people like you should be able to fight the Barbarians and win. ``The question is how, exactly, that works.''
